% Section 9: A logic generated by a satisfaction form
\label{sec:satisfaction-body}

Sections~\ref{sec:triadic-inference}--\ref{sec:through} produced carriers and crossings.
This section asks what kind of \emph{logic} such a carrier supports, and the answer is
constrained from the start: evidence is positive, propositions are sections rather than
subsets, entailment is refinement, and negations are quotients that must be declared.
The section closes by adding a quantitative enrichment and by exposing the trap that
makes it vacuous if used naively.

\subsection{Positive evidence}

\begin{definition}[Satisfaction form]\label{def:satisfaction-form}
Let $\cT_Q,\cX_Q,\cK_Q$ be the three typed carriers of a contract $Q$.  A
\emph{satisfaction form} is a relation
\[
  \Lambda_Q\subseteq \cT_Q\times\cX_Q\times\cK_Q
\]
read as a map into the Sierpinski lattice $\mathbf 2=\{\bot,\top\}$ with
$\bot<\top$, where $\top$ means ``positive evidence of satisfaction'' and $\bot$ means
``no positive evidence''.
\end{definition}

\begin{warning}[$\bot$ is not false]\label{warn:bottom-not-false}
The value $\bot$ in Definition~\ref{def:satisfaction-form} records the absence of
positive evidence.  It does not assert the negation of the proposition.  This is the
same discipline as Warning~\ref{warn:unknown-not-nonexistence}, and it is the reason the
constructions below never form a Boolean complement without declaring one.
\end{warning}

\subsection{Sections are the atoms}

\begin{definition}[Opposite-pair sections]\label{def:opposite-pair-sections}
For each pair of coordinates, define
\[
  A^{t}_{x,k}=\{\,t\in\cT_Q:\Lambda_Q(t,x,k)=\top\,\},\quad
  A^{X}_{t,k}=\{\,x:\Lambda_Q(t,x,k)=\top\,\},\quad
  A^{\cK}_{t,x}=\{\,k:\Lambda_Q(t,x,k)=\top\,\}.
\]
\end{definition}

\begin{proposal}[Generated propositions]\label{prop:generated-propositions}
The admissible propositions of the logic of $Q$ are generated by the sections of
Definition~\ref{def:opposite-pair-sections}, not by the ambient powerset:
\[
  \cL_Q=\langle G_Q\rangle,\qquad\text{and in general}\qquad
  \cL_Q\neq\cP(\cT_Q\times\cX_Q\times\cK_Q).
\]
\end{proposal}

\begin{remark}[Why this is the central design choice]
Taking the powerset as the proposition set is the default of classical semantics and it is
wrong here for a concrete reason: a subset of the carrier may be undefinable by any
section, so admitting it as a proposition silently assumes an expressive power that the
carrier does not have.  Proposal~\ref{prop:generated-propositions} refuses that
assumption and makes the expressive reach of the carrier an explicit object.
\end{remark}

\subsection{Entailment as refinement}

\begin{definition}[Support and preorder]\label{def:support-preorder}
Let $\mathrm{supp}\,P$ be the set of triples at which a generated proposition $P$ carries
value $\top$.  Define
\[
  P\preceq_Q R \iff \mathrm{supp}\,P\subseteq \mathrm{supp}\,R .
\]
\end{definition}

\begin{proposition}[Refinement is an entailment preorder]\label{prop:refinement-preorder}
$(\cL_Q,\preceq_Q)$ is a preorder, and joins exist for the opposite pairs: writing
$P_{I}$ and $P_{-I}$ for the two lifted signs of a central-sign pair,
\[
  P_{\mathrm{proj}}=P_{I}\vee P_{-I}
\]
is exactly the proposition seen by the coarse projective observer.
\end{proposition}

\begin{proof}
Reflexivity and transitivity follow from set inclusion of supports.  The join statement
is Definition~\ref{def:support-preorder} applied to the union of the two supports, which
is precisely the set on which the projective reading forgets the sign.
\end{proof}

\begin{calibration}[Observer quotient with residual]\label{cal:observer-quotient}
On the finite carrier of this section, a three-coordinate observer quotient of the
section lattice is computed together with its exact preimage table.  The residual of the
quotient is therefore known and finite, as
Proposition~\ref{prop:residual-principle} requires.
\end{calibration}

\begin{proposal}[Adjoint quantifiers]\label{prop:adjoint-quantifiers}
The quantifiers appropriate to this logic are adjoints to pullback along an observer map
$q$:
\[
  \exists_q\dashv q^{\ast}\dashv\forall_q ,
\]
where $q^{\ast}$ is the observer quotient of Calibration~\ref{cal:observer-quotient}.
\end{proposal}

\begin{remark}
Proposal~\ref{prop:adjoint-quantifiers} is a \textsf{Proposal}: the adjunction is the
right shape for observer transport, but no theorem here establishes that the finite
lattice of this section is closed under it.  What is established is the data on which
such a theorem would be stated.
\end{remark}

\subsection{The triadicity test}

\begin{proposal}[Falsifiable triadicity]\label{prop:triadicity-test}
Triadicity of $\Lambda_Q$ is falsifiable by ambiguous global fillings: if two distinct
global fillings have equal pairwise sections but different triple support, then the triple
carries information that no pair determines, and the form is triadic in the sense of
Proposal~\ref{prop:coherence-defect}.  If every two fillings with equal pairwise sections
have equal triple support, the form is only pairwise in disguise.
\end{proposal}

\begin{calibration}[The finite fixture fails the test]\label{cal:triadicity-fails}
On the finite satisfaction form constructed in this section the test of
Proposal~\ref{prop:triadicity-test} \emph{fails}: the triple support is determined by the
pairwise sections.  This is the same phenomenon as
Theorem~\ref{thm:defect-decomposes} in a logical setting, and it is the reason
Conjecture~\ref{conj:triadic-crossing} is stated as a conjecture.
\end{calibration}

\subsection{Logic as a learned characteristic}

The conceptual reading of the section is worth stating separately, because it organises
the three notions that the constructions above keep apart.

\begin{proposal}[Three grades of logicality]\label{prop:three-grades}
\begin{enumerate}
\item A \emph{proposition} is a positive (Sierpinski-valued) characteristic.
\item \emph{Entailment} is the first relational characteristic: it judges pairs.
\item \emph{Abstraction} is a coarse invariant together with a declared forgotten fibre
  $R_Q(b)$; a logic earns stability exactly by declaring what it forgets.
\end{enumerate}
\end{proposal}

\begin{warning}[Three different things]\label{warn:bivalence-completeness-decidability}
Bivalence, semantic completeness, and decidability are independent.  A logic may be
two-valued and incomplete, or complete and undecidable, and the passage from positive
evidence to a Boolean one is a \emph{second, declarable quotient} rather than a
consequence of the first.  Nothing in this section licenses forming a complement without
that second declaration.
\end{warning}

\subsection{Energy as an enrichment, and its trap}

A satisfaction form may be enriched by a quantitative measure of how a characteristic
varies along declared directions.

\begin{definition}[Directional energy]\label{def:directional-energy}
On a finite observation carrier with declared move directions $D$, the
\emph{directional Dirichlet energy} of a function $f$ is
\[
  E_D(f)=\sum_{(u,v)\in D}\bigl(f(u)-f(v)\bigr)^{2},
\]
and the \emph{energy vector} is $E_Q(f)=\bigl(E_t(f),E_X(f),E_{\cK}(f)\bigr)$.
\end{definition}

\begin{proposition}[Exact coarse--residual split]\label{prop:energy-split}
Let $q$ be an observer that averages along a forgotten fibre, and write
$\bar f$ for the coarse average and $r=f-q^{\ast}\bar f$ for the zero-mean residual of
each fibre.  Then
\[
  E_Q(f)=E_Q\bigl(q^{\ast}\bar f\bigr)+E_Q(r),
\]
and for the constructional direction in which the fibre is forgotten,
$E_{\cK}(q^{\ast}\bar f)=0$.
\end{proposition}

\begin{proof}
Each summand $\bigl(f(u)-f(v)\bigr)^{2}$ decomposes as the sum of the squares of the
coarse difference and the residual difference plus a cross term that vanishes because $r$
has zero mean on each fibre and the fibres are disjoint.  The second claim is immediate
since $q^{\ast}\bar f$ is constant along the forgotten direction.
\end{proof}

\begin{proposition}[Energy detects invariance]\label{prop:energy-invariance}
For a binary characteristic $\chi_U$,
\[
  E_D(\chi_U)=0 \iff \chi_U \text{ is constant along every } D\text{-move},
\]
and in that case $E_D$ equals the size of the boundary cut of $U$ along $D$.
\end{proposition}

\begin{proof}
Each summand is $(f(u)-f(v))^{2}\in\{0,1\}$ for binary $f$, so the sum vanishes exactly
when no edge crosses the boundary of $U$; otherwise each crossing edge contributes $1$.
\end{proof}

\begin{nogo}[The zero-energy trap]\label{nogo:zero-energy-trap}
The constant characteristic minimizes every directional energy.  Consequently any
optimisation principle of the form ``minimise $E_Q$'' is vacuous unless accompanied by an
explicit nontriviality constraint, and no learning rule may be justified by energy
decrease alone.
\end{nogo}

\begin{proof}
By Proposition~\ref{prop:energy-invariance} a constant characteristic has zero energy,
which is the global minimum since every summand is a square.
\end{proof}

\begin{remark}
No-go~\ref{nogo:zero-energy-trap} is worth stating at length because the error it forbids
is attractive: the energy splits exactly, it is computable, and its minimisers are
therefore easy to find --- too easy.  The exactness of
Proposition~\ref{prop:energy-split} and the vacuity of
No-go~\ref{nogo:zero-energy-trap} are two sides of the same computation.
\end{remark}

\subsection{Summary}

\begin{enumerate}
\item Evidence is positive and $\bot$ is not false
  (Definition~\ref{def:satisfaction-form}, Warning~\ref{warn:bottom-not-false}).
\item Propositions are generated by sections, not by the powerset, and entailment is
  refinement of support (Proposal~\ref{prop:generated-propositions},
  Proposition~\ref{prop:refinement-preorder}).
\item The triadicity test is falsifiable, and the finite fixture of this section fails it
  (Proposal~\ref{prop:triadicity-test}, Calibration~\ref{cal:triadicity-fails}).
\item Energy enriches the picture exactly, and is vacuous without a nontriviality
  constraint (Propositions~\ref{prop:energy-split}, \ref{prop:energy-invariance},
  No-go~\ref{nogo:zero-energy-trap}).
\end{enumerate}
